\subsection{Hardware validation setup}
\label{sec:method:hardware}

The simulation study deliberately omits aerodynamic drag, motor asymmetry, battery sag
and actuator wear, so its absolute errors are optimistic and only its \emph{relative}
comparisons are claimed. To test whether the ranking survives contact with a real
vehicle, the controller recommended by the benchmark is validated on hardware in a
programmable wind tunnel with motion-capture ground truth. The setup is described here;
the measurements are reported in Sec.~\ref{sec:results:hardware}.

\paragraph{Vehicle.}
A Crazyflie~2.1 Brushless --- the airframe the simulation is parameterized on
(\SI{43.4}{\gram}, thrust-to-weight $1.88$) --- carrying a Lighthouse positioning deck,
so that the onboard state estimate available in flight is the physical counterpart of
the measurement model of Sec.~\ref{sec:method:sensor}.
% TODO: state where each controller executes (onboard firmware vs. offboard at 100 Hz
% over the radio) and the payload attachment method, once the bench is assembled.

\paragraph{Programmable wind field.}
Wind is generated by a WindShaper fan-array wind
generator\footnote{Manufacturer specifications: \url{https://cdn.tytorobotics.com/Documents/Datasheets/Wind-Tunnel-Datasheet.pdf} and \url{https://www.dantecdynamics.com/solutions/fluid-mechanics/windshaper/}.} installed at the Institute of
Mechanics, Materials and Civil Engineering, UCLouvain.\footnote{\url{https://uclouvain.be/en/research-institutes/immc/news/a-new-experimental-windshaper-for-fluid-mechanics-research}}
It is an array of independently commanded \emph{wind pixels}, so the field is a
programmable function of space and time at \SI{0.1}{\second} steps, with a reported
envelope of $2$--\SI{20}{\meter\per\second} and turbulence intensity programmable over
roughly $5$--$30\,\%$.
% TODO: replace the envelope figures above with the numbers measured on the UCLouvain
% installation (module count, test-section dimensions, calibrated speed and turbulence
% range at the flight volume) before submission.
That programmability is why the facility was chosen. Both simulated wind scenarios are
defined by their \emph{structure} --- a uniform steady field, and a steady field with a
\SI{0.7}{\hertz} periodic component plus broadband turbulence --- and both are directly
expressible as WindShaper set points, so the physical and simulated conditions
correspond one-to-one and the sim-to-real gap we report is a property of the vehicle and
controller rather than of a mismatched disturbance.

\paragraph{Ground truth.}
Vehicle pose is measured by an OptiTrack PrimeX~22 motion-capture
system\footnote{Manufacturer specifications: \url{https://www.optitrack.com/cameras/primex-22}.}
at \SI{360}{\hertz} with a typical three-dimensional accuracy of
\SI{\pm0.2}{\milli\meter} --- two orders of magnitude below the tracking errors under
study and more than three times the control rate. It plays two roles. First, it supplies
the true position for exactly the metric used in simulation
($\mathrm{RMSE}_{\mathrm{3D}}$, same references, same \SI{1}{\second} warm-up
exclusion), so simulated and flown numbers are directly comparable. Second, it is the
reference against which the onboard Lighthouse estimate is differenced in flight, which
turns the sensor model of Sec.~\ref{sec:method:sensor} into a testable object: the
update rate, quasi-static bias and noise levels we simulate either reproduce against
motion capture or do not.

\paragraph{Conditions and protocol.}
The flight matrix mirrors the simulated one within the constraints of the test volume:
the same Lissajous reference at the speeds the physical envelope admits, a steady wind
and a gust profile matched to the simulated force magnitudes, and a payload step of the
same \SI{23}{\percent} of vehicle weight. Each cell is flown repeatedly so that hardware
error bars are comparable in kind to the seed statistics of
Sec.~\ref{sec:method:stats}, and the warm-up exclusion, metric definition and reference
parameterization are unchanged from simulation.
% TODO: fix the repetition count and the achievable speed tiers once the volume and the
% tunnel calibration are known.

\paragraph{What the campaign is designed to falsify.}
Three simulation claims are stated in advance as the campaign's targets, so that a
negative outcome is informative: \emph{(i)}~that under degraded onboard sensing with
wind the four leading stacks are statistically indistinguishable
(Sec.~\ref{sec:results:deployment}); \emph{(ii)}~that the hybrid adaptive stack absorbs
a payload step so completely that its tracking error is unchanged from its own nominal
value; and \emph{(iii)}~that the variance behaviour predicted under the simulated
Lighthouse model --- occasional near-failures on unlucky bias draws rather than a
uniform degradation --- appears with a real positioning deck.
